The Golden Company had wintered in worse places than King’s Landing, Ser Roland Storm reminded himself most mornings, though he was finding it harder each week to remember where, exactly, those worse places had been.

The city was cold in a way the Disputed Lands never had been — a wet, seeping cold that got into a man’s boots and stayed there, and into his mood not long after — and the gold that had drawn ten thousand sellswords across the Narrow Sea had a way of feeling considerably less warming once a man understood he’d be spending the winter garrisoned in a capital that resented his presence, fighting a war that had not yet arrived, for a queen most of the Company had never laid eyes on and half of them had begun, quietly, to doubt.

“They’re talking again,” said Black Balaq’s successor — a Summer Islander named Ossric who’d taken over the archers when the old captain finally succumbed to a lung-rot the crossing had started and the city’s damp had finished. He’d come to Roland’s quarters unbidden, which told Roland the talking had grown serious enough to warrant it. “Not just grumbling this time. Talking. Third and Fourth companies both. Wondering aloud why they’re wintering in a city that hates them to fight dead men that might not exist, when there’s a whole continent south across the water that pays just as well and doesn’t freeze a man’s piss before it hits the ground.”

“And you told them what.”

“Told them a contract’s a contract, and the Company’s never broken one in living memory, and I didn’t intend to be standing near the man who broke the first.” Ossric’s dark face held no particular conviction in it. “They didn’t much care for the answer. I confess I’m not certain I much cared for it myself, saying it. A contract’s a contract right up until the man who signed it can’t pay, and I’ve begun to wonder whether Cersei Lannister’s gold runs as deep as her promises did.”

Roland said nothing to that for a long moment, because it was precisely the question he’d been turning over himself in the small hours, and he had not yet found an answer he trusted enough to give another man.

\sceneBreak

He went to see the boy that evening, as he’d taken to doing once or twice a week — Young Griff, kept in his quiet merchant’s house well back from the Red Keep’s business, guarded and hidden and growing, Roland thought, more restless by the week beneath the enforced idleness of being a banner no one was presently marching behind.

“You’ve the look of a man who’s brought bad news,” the boy said, in that careful, schooled way of his. “You always stand straighter when it’s bad. As though good posture might make it land softer.”

“The men are restless. More than restless.” Roland did not soften it, because he’d decided a long while back that this particular boy, whatever else he was, deserved the truth plainer than most men got it. “They didn’t sign for a winter siege against something out of a fireside tale, and they signed even less for wintering idle in a city that resents them while the pay grows slower and the promises grow thinner. I’ll hold them a while yet. I don’t know how long a while is going to be, and I’d not have you surprised when it turns out to be shorter than either of us hoped.”

“And me?” The boy’s voice held something Roland had not heard in it before — not quite fear, but the nearer edge of it. “If the Company breaks, if it sails south for cleaner work, what becomes of a prince no one’s marching behind anymore?”

“That,” Roland said, honestly, “depends entirely on whether you’re actually a prince, or whether you’re a very well-trained boy Jon Connington spent a lifetime wanting to believe in. And I’ve told you before, I don’t know the answer, and I’ve stopped pretending the not-knowing troubles me the way it troubles you.” He studied the boy a moment, this handsome, unplaceable youth who might have been the last living son of Rhaegar Targaryen or might have been a Volantene bravo’s clever gamble, and found, as he always found, that he could not settle the question by looking. “I’ll tell you what I told you on the ship, since it’s the only honest counsel I have to give. The men need you to be real enough to die for. If the Company breaks, you’ll need to find something else to be real enough for, quickly, because a pretender with no army behind him is worth considerably less than a pretender with one, and considerably more dangerous to everyone who’d rather he stayed quiet.”

The boy did not answer that, and Roland did not press him, because there was nothing left to say that either of them didn’t already understand plainly enough.

\sceneBreak

Euron Greyjoy summoned him three days later, which surprised Roland, because the Crow’s Eye generally troubled himself with the Company only when he wanted something from it, and Roland had not heard of anything the man might presently want.

“Your men are unhappy,” Euron said, without preamble, standing at a window of the Red Keep that looked out over a harbor still crowded with his own quiet fleet. “Don’t trouble denying it. Half this city’s unhappy, and the other half’s too frightened to be anything at all. I find myself curious what a Golden Company captain intends to do about ten thousand unhappy sellswords, when the pay’s running thin and the war they signed for keeps declining to arrive.”

“Hold them,” Roland said. “As long as I can. It’s what a captain does.”

“And when you can’t?”

Roland met the man’s one good eye, and found in it the particular flat, amused attention that he’d learned, across a long voyage, to distrust more than open hostility. “Then I’ll do what’s best for the men under me, Lord Greyjoy. Same as I’ve always done. Whatever that turns out to be when the moment comes.”

Euron smiled at that, slow and unbothered, as though Roland had confirmed something he’d already suspected rather than threatened anything at all. “An honest answer. I’ve always liked that about you, Ser Roland — you don’t dress your loyalties up prettier than they are. Most men in your position would have sworn me some grand oath about the Company’s honor. You told me you’d look after your men and let the rest sort itself out.” He turned back to the window, to his fleet, to whatever he actually spent his hours thinking about behind that pleasant, patient face. “Keep looking after them, then. I find men who look after their own tend to end up standing on the winning side of things, when everyone grander than them has finished dying for causes that stopped mattering. I’d not mind having a man like that owe me a favor, when this is finished.”

Roland left the Red Keep with considerably more to consider than he’d arrived with, and understood, walking back down toward the quarter where his restless, unhappy men waited on a captain to tell them what came next, that Euron Greyjoy had just offered him something — not a bribe, not yet, but the shape of one, the outline of a door left deliberately ajar. He did not yet know whether he meant to walk through it. He knew only that a man who looked after his own, in a city this cold and a war this uncertain, might yet have to choose between the contract he’d sworn and the men he’d sworn it to keep alive, and that when that choice came, he intended to make it with his eyes open, whatever it cost the honor of a Company that had never, in living memory, broken a contract it had signed.